% BHU PYQ -- Manually transcribed & error-corrected
% Paper: BCH-111 : Financial Accounting - I
% Exam: B.Com. (Hons.) Semester I Examination, 2013-14

\documentclass[11pt,a4paper]{article}
\usepackage[a4paper,top=2.5cm,bottom=2.5cm,left=2.8cm,right=2.8cm,headheight=14pt]{geometry}
\usepackage{amsmath,amssymb}
\usepackage[T1]{fontenc}
\usepackage[utf8]{inputenc}
\usepackage{lmodern,microtype,fancyhdr,enumitem,hyperref,lastpage,parskip}

\hypersetup{
    colorlinks=true,
    urlcolor=black,
    linkcolor=black,
    pdftitle={BCH-111: Financial Accounting - I -- BHU B.Com. (Hons.) Sem I 2013-14}
}

\pagestyle{fancy}
\fancyhf{}
\renewcommand{\headrulewidth}{0.4pt}
\renewcommand{\footrulewidth}{0.4pt}
\fancyhead[L]{\small\textit{Banaras Hindu University}}
\fancyhead[R]{\small\textit{BCH-111: Financial Accounting - I}}
\fancyfoot[C]{\small Page \thepage\ of \pageref{LastPage}}

\newlist{parts}{enumerate}{2}
\setlist[parts,1]{label=(\roman*),leftmargin=*,itemsep=5pt,topsep=3pt}
\setlist[parts,2]{label=(\alph*),leftmargin=*,itemsep=3pt,topsep=2pt}
\newcommand{\pts}[1]{\hfill\textit{[#1]}}

\begin{document}

\begin{center}
  {\large\bfseries BANARAS HINDU UNIVERSITY}\\[2pt]
  {\normalsize B.Com. (Hons.) --- Semester I Examination, 2013-14}\\[6pt]
  \rule{0.8\linewidth}{0.5pt}\\[6pt]
  {\large\bfseries Paper No. BCH-111: Financial Accounting - I}\\[4pt]
  {\normalsize Time: 3 Hours\hfill Maximum Marks: 70}
\end{center}

\rule{\linewidth}{0.4pt}

\smallskip
\noindent\textit{Instructions: Attempt all questions. Marks are noted against each question.}

\bigskip

\noindent\textbf{Question 1.}
Define accounting. Explain the Generally Accepted Accounting Principles (GAAP) as they make accounting a perfect language of business.\pts{9+5}

\centerline{\textbf{OR}}

What is a journal? What is its use in accounting? Describe illustratively the various types of journal used in accounting.\pts{2+4+8}

\medskip\rule{\linewidth}{0.2pt}

\noindent\textbf{Question 2.}
The following errors were found in the books of a firm on December 31, 2012:
\begin{parts}
  \item Rs. 3,000 spent on the extension of factory building was charged to Repairs Account.
  \item Shri Ram's cheque of Rs. 240 was dishonoured. His account was posted in Allowance Account.
  \item Goods of Rs. 800 were returned by Ramesh Chandra. The goods were included in stock but no record for it was made in the books of account.
  \item Rs. 50 spent on the repair of a machine was posted in Machinery Account.
  \item Rs. 4,500 spent on the extension of Plant and Machinery was posted to Wages Account.
  \item Shri Ram's cheque of Rs. 150 received through Sita Ram was credited to Sita Ram's Account.
  \item An amount of Rs. 300 withdrawn by the proprietor for his personal use was posted to Travelling Expenses Account.
\end{parts}
Pass the rectifying journal entries.\pts{7}

\centerline{\textbf{OR}}

Define provision and reserve and explain the distinction between the two.\pts{7}

\medskip
\noindent\textbf{Question 2(b).}
A Company Ltd. maintains a provision for bad debts at 5\% and a provision for discount at $2\frac{1}{2}\%$ on debtors. On 01.01.2011, the balance in the provision for bad debts account was Rs. 10,000 and that in the provision for discount on debtors was Rs. 5,000. Total debtors on 31.12.2011 were Rs. 2,40,000 after writing off bad debts Rs. 6,000 and allowing discount Rs. 2,000. On 31.12.2012, total debtors were Rs. 2,00,000 after writing off bad debts Rs. 1,000 and allowing discount Rs. 500.

Show the provision for bad debts and provision for discount on debtors accounts for the two years.\pts{7}

\pagebreak

\noindent\textbf{Question 3.}
Following is the Balance Sheet of A and B on September 30, 2013. They share profits and losses in the ratio of $3/5$ and $2/5$ respectively.

\begin{center}
\small
\begin{tabular}{|l|r||l|r|}
\hline
\textbf{Liabilities} & \textbf{Amount (Rs.)} & \textbf{Assets} & \textbf{Amount (Rs.)} \\
\hline
Creditors & 15,000 & Freehold Property & 10,000 \\
Bills Payable & 3,410 & Plant and Machinery & 4,500 \\
Capitals: & & Furniture & 1,000 \\
\quad A & 18,000 & Debtors \hfill 22,500 & \\
\quad B & 16,000 & Less: Provision for Bad debts \hfill 4,000 & 18,500 \\
& & Stock & 12,500 \\
& & Investments & 4,250 \\
& & Cash & 1,660 \\
\hline
\textbf{Total} & \textbf{52,410} & \textbf{Total} & \textbf{52,410} \\
\hline
\end{tabular}
\end{center}

They admitted C into partnership from 1st October 2013. The terms of admission are as under:
\begin{parts}
  \item C to bring in Rs. 6,000 as Capital and Rs. 4,800 for goodwill in order to get $2/7$ share in profits. Both these sums are to remain in business.
  \item Goodwill of Rs. 4,800 paid by C to be credited to the Loan Account of A and B in their profit sharing ratios.
  \item Value of the whole of the goodwill on the same basis as that paid by C is to be raised in the books of the new firm.
  \item Assets are to be valued as: Freehold Property Rs. 15,000; Plant and Machinery Rs. 4,000; Stock to be discounted at 10 percent.
  \item Investments, the market value of which at the date of Balance Sheet is Rs. 3,200, to be taken over by A at that value.
  \item Provision for Bad debts is to be reduced by Rs. 1,000.
\end{parts}
C paid the sum of Rs. 10,800 by cheque on October 1, 2013.

Prepare the Revaluation Account, Capital Accounts, Loan Accounts of A and B and the Balance Sheet after C's admission.\pts{14}

\centerline{\textbf{OR}}

Define goodwill. Explain the methods of its valuation. Show the accounting treatment of goodwill on admission of a partner.\pts{14}

\medskip\rule{\linewidth}{0.2pt}

\noindent\textbf{Question 4.}
A, B and C are partners sharing profit and loss in the ratio of $3:2:2$. Their Balance Sheet as on January 1, 2012 was as follows:

\begin{center}
\small
\begin{tabular}{|l|r||l|r|}
\hline
\textbf{Liabilities} & \textbf{Amount (Rs.)} & \textbf{Assets} & \textbf{Amount (Rs.)} \\
\hline
Sundry Creditors & 16,000 & Cash & 3,000 \\
Employees' Provident Fund & 2,000 & Sundry Debtors \hfill 16,000 & \\
Reserve Fund & 4,000 & Less: Provision \hfill 1,000 & 15,000 \\
Workmen's Compensation Fund & 6,000 & Stock & 25,000 \\
Capitals: & & Furniture & 6,000 \\
\quad A \hfill 50,000 & & Patents & 4,000 \\
\quad B \hfill 30,000 & & Building & 60,000 \\
\quad C \hfill 25,000 & 1,05,000 & Goodwill & 20,000 \\
\hline
\textbf{Total} & \textbf{1,33,000} & \textbf{Total} & \textbf{1,33,000} \\
\hline
\end{tabular}
\end{center}

\pagebreak

C retires on the above date and partners agreed that:
\begin{parts}
  \item Goodwill is to be valued at two years' purchase of average profit of last four years. Profits were: 2009 -- Rs. 14,000; 2010 -- Rs. 18,000; 2011 -- Rs. 20,000 (loss); 2012 -- Rs. 15,400.
  \item Provision for doubtful debts be made at 5\% on debtors.
  \item Patents were valueless.
  \item Stock to be appreciated by 10\%.
  \item Building to be appreciated by 20\%.
  \item Furniture to be depreciated by 10\%.
  \item Creditors were to be paid Rs. 2,000 more than the book value.
  \item Goodwill shall not appear in the books of the new firm.
\end{parts}
Prepare Revaluation Account, Capital Accounts and Balance Sheet of the firm after C's retirement.\pts{14}

\centerline{\textbf{OR}}

What is a joint life policy? Show illustratively various accounting treatments of joint life policy in the books of a firm.\pts{14}

\medskip\rule{\linewidth}{0.2pt}

\noindent\textbf{Question 5.}
A, B and C are partners in a firm. Their Balance Sheet as on March 31, 2013 is given as under:

\begin{center}
\small
\begin{tabular}{|l|r||l|r|}
\hline
\textbf{Liabilities} & \textbf{Amount (Rs.)} & \textbf{Assets} & \textbf{Amount (Rs.)} \\
\hline
Creditors & 32,000 & Bank Account & 4,000 \\
Reserve & 2,000 & Machinery & 20,000 \\
Capitals: & & Stock & 8,000 \\
\quad A \hfill 15,000 & & Furniture & 2,000 \\
\quad C \hfill 6,000 & 21,000 & Debtors & 18,000 \\
& & B's Capital Account & 3,000 \\
\hline
\textbf{Total} & \textbf{55,000} & \textbf{Total} & \textbf{55,000} \\
\hline
\end{tabular}
\end{center}

The firm is dissolved due to insolvency of B who is unable to contribute anything in the payment of his dues to the firm.

Machinery realised Rs. 15,000; Furniture Rs. 3,200 and Debtors Rs. 12,000. Creditors were paid at a discount of 5\%.

Prepare the necessary accounts to close the books of the firm and show the settlement of accounts of the partners.\pts{14}

\centerline{\textbf{OR}}

Distinguish between Revaluation Account and Realisation Account. When and how are the two accounts prepared?\pts{4+2+8}

\vfill
\rule{\linewidth}{0.4pt}
\begin{center}\small
\href{https://bhu-exam.vercel.app/}{Click here to visit our website}\\[4pt]
\href{https://me-aryan.vercel.app/}{Click here to visit our contributor}\\[4pt]
\href{https://quashan-me.vercel.app/}{Click here to visit our contributor}
\end{center}

\end{document}
